\documentclass[aspectratio=169]{beamer}

% Lingua e caratteri (compilare con XeLaTeX)
\usepackage{fontspec}
\usepackage[italian]{babel}
\usepackage{amsmath}

\usetheme{default}
\setbeamertemplate{navigation symbols}{}
\setbeamertemplate{footline}[frame number]
\definecolor{accento}{HTML}{6B5CE7}
\setbeamercolor{structure}{fg=accento}

\title{Titolo della presentazione}
\subtitle{Sottotitolo}
\author{Nome Cognome}
\institute{Istituzione}
\date{\today}

\begin{document}

\begin{frame}
  \titlepage
\end{frame}

\begin{frame}{Indice}
  \tableofcontents
\end{frame}

\section{Motivazione}

\begin{frame}{Motivazione}
  \begin{itemize}
    \item<1-> Primo punto.
    \item<2-> Secondo punto, che compare dopo.
  \end{itemize}
\end{frame}

\section{Risultati}

\begin{frame}{Un risultato}
  \begin{block}{Equazione delle onde}
    \[
      \frac{\partial^2 u}{\partial t^2} = c^2 \nabla^2 u
    \]
  \end{block}
\end{frame}

\begin{frame}
  \centering\Large Grazie!
\end{frame}

\end{document}
